\section{Functors and Mutability}
\subsection{Functors, cont.}
\begin{lstlisting}
module type SetSig = SetSig
  type elem
  type t
  val empty : t
  val add : elem -> t -> t
  val mem : elem -> t -> bool
end
\end{lstlisting}
This is a functor that allows an abstract set type. We can implement it using this struct
\begin{lstlisting}
module ListSet (* : SetSig *) = struct
  type elem = int 
  type t = elem list
  let empty = []
  let add x xs = x :: xs
  let rec mem x xs =
    match xs with
    | [] -> false
    | y :: ys -> if x = y then true else mem x ys
end
\end{lstlisting}
The simple equality function is baked into OCaml, but equality is definitely not nontrivial especially from the perspective of type theory.

We can instead parameterize (as shown earlier) with an equality function.
\begin{lstlisting}
module type EqType = sig
  type t
  val equals L t -> t -> bool
end

module ListSet (T : EqType) : (SetSig with type elem = T.t) = struct
  type elem = T.t
  type t = elem list
  let empty = []
  let add x xs = x :: xs
  let rec mem x xs =
    match xs with
    | [] -> false
    | y :: ys -> if T.equals x y then true else mem x ys
\end{lstlisting}
Now to create an implementation of EqType to determine string equality.
\begin{lstlisting}
module StringEq : EqType = struct
  type t = string
  let equals s1 s2 = String.lowercase_ascii s1 = String.lowercase_ascii s2
end

module StringSet = ListSet(StringEq)
\end{lstlisting}
This is somewhat restrictive as we need to provide a new EqType every time we have an object. If we wanted to create a StringListEq, we can use functors in order to reuse our StringEq program.
\begin{lstlisting}
module ListEq (T : EqType) (* : EqType *) = struct
  type t = T.t list
  let rec equals xs ys =
    match xs, ys with 
    | [] , [] -> true
    | x :: xs' , y :: ys' -> T.equals x y && equals xs' ys'
    | _ , _ -> false
end

module StringListSet = ListSet(ListEq(StringEq))
\end{lstlisting}
In order to link the three types, we must add some extra type annotations.
\begin{lstlisting}
  module ListSet (T : EqType) : (SetSig with type elem = T.t) = struct ...
  module StringEq : (EqType with type t = string) = struct ...
  module ListEq (T : EqType) : (EqType with type t = T.t list) = struct ...
\end{lstlisting}

\subsection{Mutability}
We have this code
\begin{lstlisting}
let add_one x = x + 1

let f () =
  let x = 1 in
  let y = add_one x in
  let x = 7 in
  print_endline ("Y : " ^ string_of_int y)

let _ = f ()

\end{lstlisting}
This will output 2 as it ignores the reassignment x = 7. If we want to instead recompute x we can use reference types:
\begin{lstlisting}
let add_one (x : int ref) = !x + 1
let f () =
  let x = ref 1 in
  let y = fun _ -> add_one x in
  x := 7;
  print_endline ("Y : " ^ string_of_int (y ()))

let _ = f ()
\end{lstlisting}
The new functions have signatures ref $: a' \to a'$ ref, ! $: 'a \: ref \to 'a$, and $(:=) : 'a \: ref \to 'a \to $ unit.
\\
Included are a few excercises to learn how references work.
\begin{lstlisting}
let test () =
  let f x =
    let y = ref 10 in
    y := !y + 1;
    x + !y
  in
  let y = ref 0 in
  f 2
\end{lstlisting}
This ignores the final let y and evaluates to 13. 

\begin{lstlisting}
  let f = 
  let y = ref 0 in
    fun () ->
      y := !y + 1;
      !y
\end{lstlisting}
This counts up a value for y and increases each function call.
